\chapter*{Abstract}
\label{chap:abstract}

Automated ICD coding has improved substantially in recent years, yet many systems struggle with rare disease codes because of limited training data and the high granularity of ICD-10-CM. This thesis addresses that shortfall by combining synthetic clinical text generation, knowledge-driven validation, and memory-efficient model training to enhance both accuracy and coverage for underrepresented codes. Large language models (LLMs) are employed to produce factually diverse discharge summaries, guided by ontologies such as SNOMED CT and Orphanet. A stringent multi-phase review pipeline—encompassing rule-based checks, LLM-based fact verification, and ICD feedback loops—filters out unrealistic samples and preserves label integrity. Subsequently, these validated synthetic records augment existing training datasets for multi-label classification models, including transformer-based architectures (e.g., PLM-ICD) and synonym-aware networks. Emphasis is placed on boosting recall for rare codes, where conventional methods typically falter. Improvements in macro-F1 are observed once data scarcity is alleviated through synthetic augmentation. In addition, practical constraints are addressed by integrating approaches such as gradient checkpointing, quantization (QLoRA), and knowledge distillation, enabling resource-limited systems to fine-tune sizable language models without compromising precision. Beyond performance metrics, the thesis underscores the importance of model transparency and trustworthiness. Techniques like attention visualization and SHAP-based attributions highlight key textual fragments responsible for each predicted label, offering insights that can support clinical validation and audit processes. Error analyses further uncover any systematic biases, guiding refinements to both the data generation and classification pipelines. By merging robust augmentation, structured knowledge, and carefully optimized training routines, this work advances a scalable, interpretable ICD coding framework that narrows gaps in rare disease coverage and paves the way for more comprehensive clinical AI solutions.